%% ============================================================
%%  paper2_credit.tex
%%  "From Perpetual Contracts to Islamic Credit:
%%   The Random Stopping Time Equivalence and
%%   a Riba-Free Foundation for Sukuk, Takaful, and Credit Protection"
%%  Ahmed, Bhuyan & Islam — February 2026
%% ============================================================
\documentclass[10pt,twocolumn]{article}
\usepackage[utf8]{inputenc}
\usepackage[margin=0.75in]{geometry}
\usepackage{amsmath,amssymb,amsthm}
\usepackage{graphicx}
\usepackage{hyperref}
\usepackage[numbers]{natbib}
\usepackage{booktabs}
\usepackage{abstract}
\usepackage{xurl}
\usepackage{titlesec}
\usepackage{authblk}
\usepackage{xcolor}
\usepackage{array}
\usepackage[protrusion=true,expansion=false]{microtype}
\usepackage{enumitem}
\usepackage{caption}
\usepackage{multirow}
\usepackage{bm}

%% Section spacing
\titlespacing*{\section}{0pt}{13pt}{6pt}
\titlespacing*{\subsection}{0pt}{10pt}{4pt}
\titlespacing*{\subsubsection}{0pt}{8pt}{3pt}

%% Abstract style
\renewcommand{\abstractnamefont}{\normalfont\bfseries\large}
\renewcommand{\abstracttextfont}{\normalfont\small}

%% Theorem environments
\theoremstyle{plain}
\newtheorem{theorem}{Theorem}
\newtheorem{proposition}{Proposition}
\newtheorem{corollary}{Corollary}
\newtheorem{lemma}{Lemma}
\theoremstyle{definition}
\newtheorem{definition}{Definition}
\newtheorem{remark}{Remark}
\newtheorem{example}{Example}

%% Math operators
\DeclareMathOperator{\E}{\mathbb{E}}
\newcommand{\Qa}{Q^a}
\newcommand{\R}{\mathbb{R}}
\newcommand{\Prob}{\mathbb{P}}
\newcommand{\F}{\mathcal{F}}

\definecolor{ribared}{RGB}{180,30,30}
\definecolor{halalgreen}{RGB}{30,120,60}
\definecolor{noteblue}{RGB}{20,60,160}

\hypersetup{colorlinks=true,linkcolor=blue,citecolor=blue,urlcolor=blue}

%% ============================================================
\title{\LARGE\textbf{From Perpetual Contracts to Islamic Credit:\\
The Random Stopping Time Equivalence and\\
a Riba-Free Foundation for Sukuk,\\
Takaful, and Credit Protection}}

\author[1]{Shehzad Ahmed}
\author[2]{Rafiqul Bhuyan}
\author[3]{Rubaiyat Islam}

\affil[1]{Department of Finance, Independent University, Bangladesh\\
\small\texttt{2120922@iub.edu.bd}}
\affil[2]{Department of Finance, Alabama A\&M University, USA;
Adjunct Professor, Monarch Business School Switzerland\\
\small\texttt{rafiqul.bhuyan@aamu.edu}}
\affil[3]{Department of Software Engineering, Daffodil International University,
Bangladesh\\ \small\texttt{islam.swe@diu.edu.bd}}

\date{February 2026}
%% ============================================================

\begin{document}

\twocolumn[\begin{@twocolumnfalse}
\maketitle\thispagestyle{empty}

\begin{abstract}
\noindent
We establish a formal mathematical equivalence between two previously unconnected
pricing frameworks --- the random time representation of perpetual futures
contracts \cite{Ackerer2024} and the reduced-form credit risk model of Duffie
\& Singleton \cite{DuffieSingleton1999} --- and demonstrate that this
equivalence has profound implications for Islamic finance.

The core observation is structural: the stopping time $\theta_t$ in the
Ackerer-Hugonnier-Jermann perpetual pricing formula is isomorphic to the
default time $\tau$ in the Duffie-Singleton defaultable bond formula,
under the identification of the convergence intensity $\kappa$ with the hazard
rate $\lambda$ and the interest parameter $\iota$ with the risk-free
discount rate $r$. When the interest parameter is set to zero ($\iota = 0$)
--- which Ackerer et al.\ prove achieves unique no-arbitrage
perpetual futures pricing --- the analogous credit instrument at $r = 0$
is also no-arbitrage, with the stopping time distribution alone carrying
all timing information previously encoded in the interest discount.

This equivalence resolves a 50-year impasse in Islamic finance: it provides
a mathematically rigorous, no-arbitrage pricing foundation for credit
instruments --- sukuk, takaful, and credit default swaps --- that requires
no positive interest rate. We prove the main theorem, derive closed-form
prices for an Islamic perpetual sukuk and an everlasting takaful contract
using Proposition~6 of Ackerer et al., and show that the convergence
intensity $\kappa$ functions as a riba-free analog to the risk-free rate
across all Islamic credit instruments. We term this the \textbf{$\kappa$-rate}.

\medskip
\noindent\textbf{Keywords:} \textit{riba}, Islamic finance, credit risk,
sukuk, takaful, perpetual contracts, stopping time, no-arbitrage,
Duffie-Singleton, $\kappa$-rate, riba-free pricing
\end{abstract}
\bigskip
\end{@twocolumnfalse}]

%% ============================================================
\section{Introduction}
%% ============================================================

The pricing of credit risk is the central unresolved problem in Islamic
finance. The \$3 trillion Islamic finance industry \cite{IFSB2024} operates
without any mathematical tool for pricing credit default risk that does not
require a positive interest rate. Every credit derivatives framework in use
today --- the structural models of Merton \cite{Merton1974}, Black \&
Cox \cite{BlackCox1976}, and Leland \cite{Leland1994}; the reduced-form
models of Jarrow \& Turnbull \cite{JarrowTurnbull1995}, Duffie \&
Singleton \cite{DuffieSingleton1999}, and Lando \cite{Lando1998} --- is
built on a positive risk-free interest rate $r > 0$. Islamic scholars
correctly identify this $r$ as \textit{riba}: it is the compensation for
the pure time value of money, which Islamic jurisprudence prohibits as a
predetermined excess in exchange transactions \cite{ElGamal2006,Usmani2002}.

The consequence has been severe. Islamic institutions cannot hedge credit
risk. Sukuk (Islamic bonds) are priced at issuance using LIBOR/SOFR-anchored
models, making the interest rate implicit even when explicit interest
payments are restructured away \cite{KhanAbdallah2017}. Takaful (Islamic
insurance) operators invest pools in sukuk priced by riba-based models
\cite{Htay2012}. The Islamic credit default swap does not exist. Scholars
have sought workarounds for decades; the workarounds have been cosmetic,
not mathematical \cite{Khan2010,ElGamal2006}.

This paper resolves the impasse. We demonstrate that a recent result from
mathematical finance --- the random time representation theorem for perpetual
futures contracts proved by Ackerer, Hugonnier \& Jermann (2024)
\cite{Ackerer2024} --- contains, as a special case, a complete and
rigorous pricing framework for credit instruments that does not require
$r > 0$.

\subsection*{The Core Observation}

The Ackerer-Hugonnier-Jermann theorem establishes that the unique
no-arbitrage price of a perpetual futures contract on a spot asset $X_t$
with interest parameter $\iota$ and convergence intensity $\kappa$ is:
\begin{equation}
f_t = \E^{\Qa}_t\!\left[\mathrm{e}^{-\int_t^{\theta_t}\!\iota_s\,ds}\,
X_{\theta_t}\right]
\label{eq:ahj}
\end{equation}
where $\theta_t$ is a random stopping time with exponential density
under the pricing measure $\Qa$. The critical result, proved in their
Theorem~3, is that $\iota = 0$ satisfies all no-arbitrage conditions.
When $\iota = 0$, equation~\eqref{eq:ahj} collapses to:
\begin{equation}
f_t = \E^{\Qa}_t[X_{\theta_t}], \qquad \theta_t\sim t + \mathrm{Exp}(\kappa)
\label{eq:ahj_zero}
\end{equation}
The instrument price is the \emph{expected spot value at a random future
moment} --- with the probability distribution of that moment
determined entirely by $\kappa$, and no interest discount anywhere.

The Duffie-Singleton reduced-form credit model prices a defaultable
instrument as:
\begin{equation}
P_t = \E^Q_t\!\left[\mathrm{e}^{-\int_t^{\tau}r_s\,ds}\,R_\tau\right]
\label{eq:ds}
\end{equation}
where $\tau$ is the default time (exponentially distributed with hazard
rate $\lambda$) and $R_\tau$ is the recovery value at default.

Equations~\eqref{eq:ahj} and~\eqref{eq:ds} are the same equation.
The stopping time $\theta_t$ is $\tau$. The convergence intensity
$\kappa$ is $\lambda$. The interest parameter $\iota$ is $r$.
The spot $X_{\theta_t}$ is the recovery $R_\tau$.

At $\iota = r = 0$: both yield $\E[X_\tau]$. The Ackerer theorem
proves this is no-arbitrage. Therefore credit pricing at $r = 0$ is
no-arbitrage, provided $\kappa = \lambda > 0$ carries all timing
information. \textit{Riba} is not a pricing necessity. It is a
historical artifact of how the two literatures fused timing and
discounting into a single parameter.

\subsection*{Contributions}

This paper makes five contributions.

\begin{enumerate}[leftmargin=*, noitemsep]
\item \textbf{Formal equivalence theorem.} We prove that the
Ackerer-Hugonnier-Jermann perpetual pricing formula and the
Duffie-Singleton defaultable bond formula are formally isomorphic
(Proposition~\ref{prop:equiv}), and that no-arbitrage at $\iota = 0$
implies no-arbitrage for credit pricing at $r = 0$
(Corollary~\ref{cor:riba_free}).

\item \textbf{Islamic perpetual sukuk.} We derive a closed-form
no-arbitrage price for an Islamic perpetual sukuk using the
$\kappa$-parameterized stopping time (Section~\ref{sec:sukuk}).

\item \textbf{Actuarially fair takaful formula.} We apply Proposition~6
of Ackerer et al.\ to derive a closed-form everlasting takaful
contribution pricing formula --- riba-free, on-chain verifiable, and
Shariah-compatible (Section~\ref{sec:takaful}).

\item \textbf{The $\kappa$-rate.} We propose the convergence intensity
$\kappa$ as an Islamic alternative to the risk-free rate $r$, providing
the first mathematically rigorous pricing foundation for a complete
riba-free credit market (Section~\ref{sec:kappa_rate}).

\item \textbf{Shariah analysis.} We conduct a formal jurisprudential
analysis showing that the $\kappa$-parameterized framework satisfies
the riba, gharar, and maysir conditions of AAOIFI Shariah Standards
\cite{AAOIFI2017} (Section~\ref{sec:shariah}).
\end{enumerate}

%% ============================================================
\section{The Ackerer-Hugonnier-Jermann Framework}
\label{sec:ahj}
%% ============================================================

\subsection{Market Setup}

Let $(\Omega, \F, \Prob)$ be a probability space with filtration
$(\F_t)_{t\geq 0}$ satisfying the usual conditions. There are two
funding accounts: a long account growing at rate $r_a(t)$ and a short
account growing at rate $r_b(t)$, both strictly positive. Let $X_t$
denote the price of a spot asset satisfying, under the long's pricing
measure $\Qa$:
\begin{equation}
\frac{dX_t}{X_t} = (r_a(t) - r_b(t))\,dt + \sigma_X\,dW_t^{\Qa}
\label{eq:spot}
\end{equation}
where $W^{\Qa}$ is a standard Brownian motion under $\Qa$ and
$\sigma_X > 0$ is the spot volatility. A perpetual futures contract
on $X_t$ with convergence parameter $\kappa > 0$ and interest
parameter $\iota \geq 0$ is a continuously-settled contract with
no fixed expiry.

\subsection{The Perpetual Futures Price}

Under constant parameters $r_a, r_b, \kappa, \iota, \sigma_X$,
Proposition~3 of Ackerer et al.\ \cite{Ackerer2024} establishes that
the unique no-arbitrage perpetual futures price is:
\begin{equation}
f_t = \frac{\kappa}{\kappa - \iota + r_b - r_a} \cdot X_t
\label{eq:perp_price}
\end{equation}
provided $\kappa - \iota + r_b - r_a > 0$.
The ratio $f_t/X_t$ is the perpetual \textit{basis}. When
$r_a = r_b$ (symmetric funding rates):
\begin{equation}
f_t = \frac{\kappa}{\kappa - \iota} \cdot X_t
\label{eq:perp_symmetric}
\end{equation}

\subsection{The Random Time Representation}
\label{sec:rtr}

The central result of Ackerer et al.\ is a \textit{random time
representation} that reveals the economic structure of perpetual
pricing. Define the \textbf{convergence stopping time} $\theta_t$
as a random variable with conditional density under $\Qa$:
\begin{equation}
q(\theta_t \in ds \mid \F_t)
= (\kappa_s - \iota_s)\,
\exp\!\left(-\int_t^s (\kappa_u - \iota_u)\,du\right)ds, \quad s\geq t
\label{eq:density}
\end{equation}
This is a Cox process (doubly stochastic Poisson process) with
intensity $\kappa_t - \iota_t$. The distribution of $\theta_t$
is exponential with rate $\kappa - \iota$ in the constant-parameter
case:
\begin{equation}
\Qa(\theta_t > s \mid \F_t)
= e^{-(\kappa-\iota)(s-t)}, \quad s\geq t
\label{eq:survival}
\end{equation}
The perpetual futures price then satisfies (Theorem~6, Ackerer et al.):
\begin{equation}
f_t = \E^{\Qa}_t\!\left[
\mathrm{e}^{-\int_t^{\theta_t}\iota_s\,ds} \cdot X_{\theta_t}
\right]
\label{eq:rtr_full}
\end{equation}
The price is the discounted expected spot value at the random
stopping time $\theta_t$. The discount is $\iota$ (the interest
parameter), not $r_a$ or $r_b$.

\subsection{The Critical Case: $\iota = 0$}
\label{sec:iota_zero}

When $\iota = 0$, Theorem~3 of Ackerer et al.\ establishes that
no-arbitrage is satisfied. Substituting into
equations~\eqref{eq:density}--\eqref{eq:rtr_full}:
\begin{align}
q(\theta_t \in ds \mid \F_t) &= \kappa\,e^{-\kappa(s-t)}\,ds
\label{eq:density_zero}\\[4pt]
f_t &= \E^{\Qa}_t[X_{\theta_t}],
\quad \theta_t \sim t + \mathrm{Exp}(\kappa)
\label{eq:perp_iota_zero}
\end{align}
Three consequences are immediate and profound:

\begin{enumerate}[leftmargin=*,noitemsep]
\item \textbf{No discounting.} The expectation in
\eqref{eq:perp_iota_zero} is taken without any interest discount.
The price is the pure expected value.

\item \textbf{$\kappa$ carries all timing information.} The
probability distribution of \textit{when} payoff occurs is
governed entirely by $\kappa$. Larger $\kappa$ means faster
convergence and shorter expected stopping time;
smaller $\kappa$ means slower convergence.

\item \textbf{The price formula at $r_a = r_b$.}
From \eqref{eq:perp_symmetric} with $\iota = 0$: $f_t = X_t$.
The perpetual futures price equals the spot price identically ---
no premium, no discount --- consistent with the empirical
finding that dYdX~v4 perpetuals at $\iota=0$ exhibit zero
long-run mean funding rates \cite{Ahmed2026}.
\end{enumerate}

\begin{remark}
The CEX interest parameter $\iota = 0.0001$ per 8 hours
($= 0.10950\,\%$/year on Binance, Bybit, OKX) creates a permanent
positive bias in equation~\eqref{eq:perp_symmetric}: $f_t > X_t$
always. Longs pay this excess to shorts regardless of market
conditions. This excess is structurally \textit{riba al-fadl}
(predetermined excess in exchange) under all four Sunni schools.
Setting $\iota = 0$ eliminates it exactly.
\end{remark}

%% ============================================================
\section{Reduced-Form Credit Risk Models}
\label{sec:credit}
%% ============================================================

\subsection{The Duffie-Singleton Framework}

The reduced-form approach to credit risk, initiated by Jarrow \&
Turnbull \cite{JarrowTurnbull1995} and systematized by Duffie \&
Singleton \cite{DuffieSingleton1999}, models default as the first
jump of a Poisson process with stochastic intensity (hazard rate).

Let $(\Omega, \F, Q)$ be a risk-neutral probability space. Define
the \textbf{default time} $\tau$ as a stopping time with respect
to $(\F_t)$, characterized by the $Q$-conditional survival
probability:
\begin{equation}
Q(\tau > s \mid \F_t)
= \exp\!\left(-\int_t^s \lambda_u\,du\right), \quad s\geq t
\label{eq:ds_survival}
\end{equation}
where $\lambda_t > 0$ is the \textbf{hazard rate} (default
intensity). The hazard rate has the conditional density:
\begin{equation}
q(\tau \in ds \mid \F_t)
= \lambda_s \exp\!\left(-\int_t^s \lambda_u\,du\right)ds,
\quad s\geq t
\label{eq:ds_density}
\end{equation}

\subsection{Defaultable Instrument Pricing}

For a defaultable instrument with payoff function $\phi(X_\tau)$
at the default time $\tau$ and risk-free discount rate $r_t$:
\begin{equation}
P_t = \E^Q_t\!\left[
\mathrm{e}^{-\int_t^\tau r_s\,ds} \cdot \phi(X_\tau)
\right]
\label{eq:ds_price}
\end{equation}
In the constant-parameter case ($r, \lambda$ constant, $\tau\sim
t+\mathrm{Exp}(\lambda)$), this gives:
\begin{equation}
P_t = \int_t^\infty e^{-r(s-t)}\,\phi\!\left(\E[X_s]\right)\,
\lambda e^{-\lambda(s-t)}\,ds
= \frac{\lambda}{r+\lambda}\,\phi(X_t^*)
\label{eq:ds_closed}
\end{equation}
where $X_t^*$ denotes the forward spot expectation. For a bond
with full recovery ($\phi(x) = x$) and $Q$-martingale spot:
\begin{equation}
P_t = \frac{\lambda}{r+\lambda}\cdot X_t
\label{eq:ds_bond}
\end{equation}
The conventional framework requires $r > 0$ to obtain
$P_t < X_t$: without interest discounting, the bond price
equals the spot, which appears to eliminate the credit spread.
We resolve this apparent paradox in Section~\ref{sec:equiv}.

%% ============================================================
\section{The Mathematical Equivalence}
\label{sec:equiv}
%% ============================================================

\subsection{The Isomorphism}

Comparing equations~\eqref{eq:density} and \eqref{eq:ds_density}
reveals an immediate structural identity:

\begin{center}
\small
\begin{tabular}{@{}lll@{}}
\toprule
\textbf{AHJ Perpetuals} & \textbf{Duffie-Singleton Credit} & \textbf{Role} \\
\midrule
$\theta_t$ & $\tau$ & Random stopping time \\
$\kappa_t$ & $\lambda_t$ & Event intensity \\
$\iota_t$ & $r_t$ & Discount parameter \\
$X_{\theta_t}$ & $\phi(X_\tau)$ & Payoff at event \\
$f_t$ & $P_t$ & Instrument price \\
$\Qa$ & $Q$ & Pricing measure \\
\bottomrule
\end{tabular}
\end{center}

The two frameworks share an identical mathematical skeleton. The
perpetual futures framework uses $\kappa$ and $\iota$; the credit
framework uses $\lambda$ and $r$. The functions they play are
identical.

\subsection{The Main Result}
\label{sec:main_theorem}

\begin{proposition}[Random Stopping Time Equivalence]
\label{prop:equiv}
Let $(\Omega,\F,\Prob)$ be a filtered probability space with the
usual conditions. Let $X_t$ be an asset price process satisfying
\eqref{eq:spot} under a pricing measure $Q = \Qa$. Define:
\begin{enumerate}[leftmargin=*,noitemsep]
\item A perpetual futures contract with parameters $\kappa, \iota$,
priced by equation~\eqref{eq:rtr_full}.
\item A defaultable financial claim with hazard rate $\lambda$,
discount rate $r$, and recovery $\phi(X_\tau)$, priced by
equation~\eqref{eq:ds_price}.
\end{enumerate}
Under the parameter identification $\kappa \leftrightarrow \lambda$,
$\iota \leftrightarrow r$, $X_{\theta_t} \leftrightarrow \phi(X_\tau)$:
\begin{equation}
f_t = P_t
\label{eq:equiv}
\end{equation}
The two pricing formulas are formally isomorphic.
\end{proposition}

\begin{proof}
The stopping time densities \eqref{eq:density} and
\eqref{eq:ds_density} are identical under $\kappa \leftrightarrow
\lambda$. The integral representations \eqref{eq:rtr_full} and
\eqref{eq:ds_price} are identical under the additional identification
$\iota \leftrightarrow r$ and $X_{\theta_t} \leftrightarrow \phi(X_\tau)$.
The pricing measures $\Qa$ and $Q$ are both risk-neutral with respect
to their respective numeraires. The formal identity \eqref{eq:equiv}
follows directly.\qquad$\square$
\end{proof}

\begin{corollary}[Riba-Free Credit Pricing is No-Arbitrage]
\label{cor:riba_free}
Setting $\iota = 0$ in the AHJ perpetual futures framework yields
a no-arbitrage price by Theorem~3 of \textup{\cite{Ackerer2024}}.
Under the isomorphism of Proposition~\ref{prop:equiv}, the
corresponding credit instrument at $r = 0$ is also no-arbitrage,
with pricing formula:
\begin{equation}
\boxed{P_t = \E^Q_t[\phi(X_{\tau})], \qquad
\tau\sim t+\mathrm{Exp}(\kappa), \quad r = 0}
\label{eq:riba_free}
\end{equation}
The hazard rate $\kappa$ alone carries all timing information.
No positive interest rate is required.
\end{corollary}

\begin{proof}
By Theorem~3 of Ackerer et al.\ \cite{Ackerer2024}, setting
$\iota = 0$ satisfies the no-arbitrage conditions for the perpetual
futures. By Proposition~\ref{prop:equiv}, $f_t = P_t$ under the
given identification. Therefore $P_t$ at $r = 0$ satisfies the
same no-arbitrage conditions.
Specifically: there exists a self-financing replication strategy
for the claim $\phi(X_\tau)$ with value process $P_t$ satisfying
$P_t = \E^Q_t[\phi(X_\tau)]$, as the expectation under the
risk-neutral measure $Q$ of the discounted (at $r=0$) payoff.
\qquad$\square$
\end{proof}

\subsection{Resolving the Apparent Paradox}

Corollary~\ref{cor:riba_free} may appear paradoxical: how can a
credit instrument with $r = 0$ have $P_t < X_t$ (i.e., trade at
a discount) without interest? The resolution is that $\kappa$
controls both the discount and the credit spread simultaneously.

From equation~\eqref{eq:perp_symmetric} with $\iota = 0$:
\begin{equation}
\frac{f_t}{X_t} = 1
\qquad (r_a = r_b,\;\iota = 0)
\label{eq:parity}
\end{equation}
But consider a sukuk with recovery $\phi(X_\tau) = \delta \cdot X_\tau$
where $\delta \in (0,1)$ is the recovery fraction. Then:
\begin{equation}
P_t = \delta\cdot\E^Q_t[X_\tau] = \delta\cdot X_t < X_t
\label{eq:recovery_discount}
\end{equation}
The discount from par is $1-\delta$: the credit loss on default.
This is not interest; it is the pricing of credit risk through
the recovery rate. The $\kappa$-parameterized timing generates
the spread (the probability-weighted present value of credit loss)
without any interest component. This is \textit{precisely} the
separation of credit risk from time preference that Islamic
finance requires.

\begin{remark}
In the conventional framework, the credit spread over the risk-free
rate bundles two economically distinct risks: (1) the probability
of default (hazard rate $\lambda$) and (2) the time value of money
(risk-free rate $r$). The Ackerer framework separates them
definitively. At $\iota = r = 0$, only the hazard rate $\kappa$
remains, representing pure credit risk with no interest component.
\end{remark}

%% ============================================================
\section{Everlasting Options and the Takaful Pricing Formula}
\label{sec:everlasting}
%% ============================================================

Proposition~6 of Ackerer et al.\ \cite{Ackerer2024} extends the
random time framework to arbitrary payoff functions $\phi(X_\tau)$,
yielding closed-form \textbf{everlasting option} prices. This
extension is the key to takaful pricing.

\subsection{The Everlasting Put Option}

For a payoff $\phi(x) = \max(K - x, 0)$ (floor protection at
strike $K$) with $\iota = 0$ and $r_a = r_b$, the price is:
\begin{equation}
\Pi_{\mathrm{put}}(x, K)
= \E^{\Qa}\!\left[(K - X_\tau)^+\right]
= C_p \cdot x^{\beta_p}
\label{eq:evput}
\end{equation}
where $\beta_p < 0$ is the \textit{negative} root of the
characteristic equation:
\begin{equation}
\frac{\sigma_X^2}{2}\,\beta(\beta-1) - \kappa = 0
\qquad (\iota = 0,\;r_a = r_b)
\label{eq:char_eq}
\end{equation}
Solving:
\begin{equation}
\beta_p = \frac{1}{2} - \sqrt{\frac{1}{4} + \frac{2\kappa}{\sigma_X^2}}
\label{eq:beta_p}
\end{equation}
The constant $C_p$ is determined by value-matching at $x = K$
(smooth-pasting condition):
\begin{equation}
C_p = \frac{K^{1-\beta_p}}{1 - \beta_p}
\label{eq:Cp}
\end{equation}

\subsection{The Everlasting Call Option}

For $\phi(x) = \max(x - K, 0)$, the everlasting call is:
\begin{equation}
\Pi_{\mathrm{call}}(x, K) = C_c \cdot x^{\beta_c}
\label{eq:evcall}
\end{equation}
where $\beta_c > 1$ is the positive root of \eqref{eq:char_eq}:
\begin{equation}
\beta_c = \frac{1}{2} + \sqrt{\frac{1}{4} + \frac{2\kappa}{\sigma_X^2}}
\label{eq:beta_c}
\end{equation}
with $C_c = K^{1-\beta_c}/(\beta_c - 1)$.

Note that $|\beta_p| + \beta_c = 2\beta_c - 1$, and both satisfy
put-call parity at $r=0$: $\Pi_{\mathrm{call}} - \Pi_{\mathrm{put}}
= X_t - K$ (forward parity without discounting).

%% ============================================================
\section{Applications in Islamic Finance}
\label{sec:applications}
%% ============================================================

\subsection{Islamic Perpetual Sukuk}
\label{sec:sukuk}

\begin{definition}[Islamic Perpetual Sukuk]
An Islamic Perpetual Sukuk is a financial instrument in which:
\begin{enumerate}[leftmargin=*,noitemsep]
\item The sukuk holder contributes capital $K$ backed by a real
asset (RWA) with market value $X_t$.
\item The issuer (asset operator) pays a continuous premium
$\pi_t = \kappa(f_t - X_t)\,dt$ to the holder --- the perpetual
futures funding rate under the $\kappa$-parameterized formula.
\item At a random redemption time $\tau$ (oracle-triggered by
asset completion, sale, or Shariah board review), the holder
receives the fair market value $X_\tau$ of the underlying asset.
\item The interest parameter $\iota = 0$ at all times.
\end{enumerate}
\end{definition}

\begin{proposition}[Sukuk No-Arbitrage Price]
\label{prop:sukuk}
The fair price of the Islamic Perpetual Sukuk at time $t$ is:
\begin{equation}
P_t^{\mathrm{sukuk}} = \E^Q_t[X_\tau]
= \frac{\kappa}{\kappa + r_b - r_a}\cdot X_t
\label{eq:sukuk_price}
\end{equation}
with $r_a = r_b$ (symmetric bilateral rates) giving $P_t = X_t$.
This price is no-arbitrage by Corollary~\textup{\ref{cor:riba_free}}.
\end{proposition}

\begin{proof}
The redemption time $\tau \sim t + \mathrm{Exp}(\kappa)$ under
the pricing measure $Q$. With recovery $\phi(X_\tau) = X_\tau$,
Corollary~\ref{cor:riba_free} gives $P_t = \E^Q_t[X_\tau]$.
Under the geometric Brownian motion \eqref{eq:spot}:
$\E^Q_t[X_\tau] = \E^Q_t[X_t\,e^{(r_a-r_b)(\tau-t)}] =
X_t \cdot \kappa/(\kappa + r_b - r_a)$, which is equation
\eqref{eq:sukuk_price}. No-arbitrage follows from
Corollary~\ref{cor:riba_free}.\qquad$\square$
\end{proof}

\paragraph{Economic interpretation.}
The sukuk is priced at the expected asset value at redemption,
with no interest discount. The premium payments $\pi_t$ transfer
the basis risk from the issuer to the holder: when the asset
trades at a premium to spot ($f_t > X_t$), the operator pays
the holder; when at a discount, the holder pays the operator.
This is risk-sharing (\textit{mudarabah}-compatible) rather than
a fixed return on capital.

\paragraph{The $\kappa$ parameter for sukuk.}
For a sukuk backed by an infrastructure asset, $\kappa$ is
calibrated from the expected completion timeline. An asset with
a 5-year expected completion horizon has $\kappa \approx 0.20$
(rate of an exponential with mean 5 years). A 10-year asset:
$\kappa \approx 0.10$. This is directly observable and
estimable without any interest rate input.

\subsection{Takaful: Actuarially Fair Perpetual Insurance}
\label{sec:takaful}

\begin{definition}[Everlasting Takaful Contract]
An Everlasting Takaful Contract is a mutual insurance arrangement
in which:
\begin{enumerate}[leftmargin=*,noitemsep]
\item Participants contribute premium $\pi_\mathrm{takaful}$
into a mutual pool at inception.
\item The pool pays a claim $\phi(X_\tau) = \max(K - X_\tau, 0)$
--- the shortfall of asset value below the insured floor $K$
--- at a random claim event $\tau$.
\item The claim event follows a Cox process with intensity
$\kappa$ (the historical hazard rate of the insured event:
drought, flood, death, property damage).
\item No investment return is generated on the pool; it earns
at rate $\iota = 0$.
\end{enumerate}
\end{definition}

\begin{theorem}[Actuarially Fair Takaful Premium]
\label{thm:takaful}
The actuarially fair takaful contribution for floor protection
at level $K$ when the current asset value is $x = X_t$, event
intensity $\kappa$, and asset volatility $\sigma_X$ is:
\begin{equation}
\boxed{\Pi_\mathrm{takaful}(x, K)
= \frac{K^{1-\beta_p}}{1-\beta_p}\cdot x^{\beta_p}}
\label{eq:takaful_premium}
\end{equation}
where $\beta_p = \frac{1}{2} - \sqrt{\frac{1}{4}
+ \frac{2\kappa}{\sigma_X^2}} < 0$.
This price is no-arbitrage (Corollary~\textup{\ref{cor:riba_free}})
and requires no positive interest rate.
\end{theorem}

\begin{proof}
By Corollary~\ref{cor:riba_free}, the no-arbitrage price of
$\phi(X_\tau) = (K-X_\tau)^+$ at $r=0$ is $\E^Q_t[(K-X_\tau)^+]$.
By Proposition~6 of Ackerer et al.\ \cite{Ackerer2024} applied
with $\iota = 0$ and $r_a = r_b$, this expectation equals
$C_p x^{\beta_p}$ with $C_p$ and $\beta_p$ as given in
equations~\eqref{eq:Cp} and~\eqref{eq:beta_p}. No-arbitrage
is inherited from Corollary~\ref{cor:riba_free}.\qquad$\square$
\end{proof}

\paragraph{Agricultural takaful example.}
Let $X_t$ be the market price of wheat (USD/bushel). A farmer
participates in a mutual takaful pool offering floor protection
at $K = \$6.00$/bushel. Historical drought frequency in the
region is $\kappa = 0.08$ (one drought every $\sim$12.5 years
on average). Wheat price volatility is $\sigma_X = 0.25$
(25\% annual). Current wheat price: $x = \$7.50$.

\begin{align}
\beta_p &= \frac{1}{2} - \sqrt{\frac{1}{4}
+ \frac{2(0.08)}{(0.25)^2}}
= 0.5 - \sqrt{0.25 + 2.56} = 0.5 - 1.676 = -1.176
\notag\\[4pt]
C_p &= \frac{(6.00)^{1-(-1.176)}}{1-(-1.176)}
= \frac{6.00^{2.176}}{2.176}
= \frac{47.85}{2.176} = 21.99
\notag\\[4pt]
\Pi &= 21.99 \times (7.50)^{-1.176}
= 21.99 \times 0.1058 = \$2.33\text{/bushel}
\label{eq:agri_example}
\end{align}

The farmer contributes \$2.33 per bushel into the takaful pool.
This is the actuarially fair premium. No interest rate was used.
The contribution is fully transparent and on-chain verifiable.

\paragraph{Life takaful.}
Let $X_t$ be a human capital index (proxied by a broad equity
index). Let $\kappa$ be the age-adjusted mortality hazard rate
(from actuarial tables). The death benefit is $\phi(X_\tau) = B$
(a fixed sum). From Proposition~6 with $\phi(x) = B$ (constant),
$\Pi_\mathrm{life} = B \cdot \E^Q[e^{-\kappa(\tau-t)}] =
B \cdot \kappa/\kappa = B$ evaluated at the actuarially discounted
rate. More precisely, the fair contribution for a benefit $B$
payable at death $\tau \sim \mathrm{Exp}(\kappa)$ is:
$\Pi = B \cdot \kappa / (\kappa + 0) = B$, which with realistic
$B < X_t$ gives $\Pi < X_t$ --- the participant contributes
less than the death benefit at any point in time, with the
pool providing mutual guarantee (\textit{takaful}).

\subsection{Islamic Credit Default Swaps (iCDS)}
\label{sec:icds}

The global credit default swap (CDS) market has a notional
outstanding of approximately \$10 trillion \cite{BIS2024}. It
is entirely built on the Duffie-Singleton framework with $r > 0$.
Islamic institutions have no equivalent instrument for credit
risk hedging.

\begin{definition}[Islamic CDS (iCDS)]
An iCDS is a bilateral contract in which:
\begin{enumerate}[leftmargin=*,noitemsep]
\item The protection buyer pays a continuous premium $\kappa\,dt$
per unit time to the protection seller.
\item At the credit event time $\tau$ (default, restructuring, or
rating trigger), the seller pays $X_\tau$ (the fair market value
of the reference sukuk at that moment) to the buyer.
\item The interest parameter $\iota = 0$; no fixed rate premium.
\end{enumerate}
\end{definition}

\begin{proposition}[iCDS Pricing]
The fair value of the iCDS protection to the buyer at time $t$
when the reference sukuk has current value $X_t$, default
intensity $\kappa$, and $\iota = 0$ is:
\begin{equation}
V_t^{\mathrm{iCDS}} = \E^Q_t[X_\tau]
- \kappa\int_t^\infty e^{-\kappa(s-t)}\,ds
= X_t - 1
\label{eq:icds}
\end{equation}
for normalized $X_t$ (sukuk price in units of par). The fair
$\kappa$ (iCDS spread) is implied from the observed market
basis of the reference sukuk.
\end{proposition}

The iCDS premium $\kappa$ is calibrated from market observables
(the sukuk basis, historical default rates of the reference
entity, rating agency transition matrices) --- exactly as
conventional CDS spreads are calibrated from bond yield spreads,
with no interest rate input required.

\subsection{Sovereign Perpetual Sukuk}
\label{sec:sovereign}

Sovereign governments in Muslim-majority countries (Malaysia,
Saudi Arabia, Indonesia, Türkiye, Pakistan) issue sukuk to fund
infrastructure. These are priced at issuance against SOFR or
government bond yields, embedding riba implicitly
\cite{IIFM2023,IFSB2024}. The Islamic perpetual sukuk framework
resolves this.

\begin{example}[Infrastructure Perpetual Sukuk]
Consider a sovereign issuing a perpetual sukuk backed by a
highway project with construction timeline of 7 years (expected
completion, exponentially distributed): $\kappa = 1/7 \approx
0.143$. The project asset value is $X_t$ (updated quarterly
by a government oracle). The sukuk trades at $P_t = X_t$
(at-par perpetual under $r_a = r_b$, $\iota = 0$).

The sovereign pays the continuous funding premium:
$\pi_t = \kappa(f_t - X_t)\,dt$ to sukuk holders.
On completion (random event $\tau$), bondholders receive $X_\tau$.

Secondary market price at any time: $P_t = X_t$ (adjusted
for observed $\kappa$ from reported completion timeline).
\emph{No yield-to-maturity calculation required. No SOFR
benchmark needed. No riba.}
\end{example}

Critically, the $\kappa$ parameter serves as a \textit{sovereign
credibility signal}: a government managing its infrastructure
responsibly exhibits stable, predictable $\kappa$; governance
failures manifest as declining implied $\kappa$ (basis widening),
creating a market-based credit signal without any interest rate.

%% ============================================================
\section{The $\kappa$-Rate: An Islamic Monetary Alternative}
\label{sec:kappa_rate}
%% ============================================================

\subsection{The Problem with $r$}

The risk-free interest rate $r$ is the fundamental pricing
parameter of conventional finance. It is:
\begin{enumerate}[leftmargin=*,noitemsep]
\item Set by central bank monetary policy (exogenous, not
determined by production or real economic activity).
\item Defined as the return to lending money at zero credit risk
--- the pure reward for the time value of money.
\item Definitionally \textit{riba}: it compensates capital for
time elapsed, regardless of real economic outcome.
\end{enumerate}
Every discount factor in every asset pricing model contains $r$.
The equity risk premium, the bond yield, the option delta ---
all are expressed relative to $r$. Islamic finance has no
alternative. Practitioners either (i) use $r$ despite the
scholarly objection, or (ii) use LIBOR/SOFR as a \textit{benchmark}
while restructuring cash flows, which is cosmetic compliance.

\subsection{The $\kappa$-Rate as a Substitute}

We propose the \textbf{$\kappa$-rate}: the convergence intensity
parameter of the Ackerer-Hugonnier-Jermann framework as a
riba-free substitute for $r$ in Islamic finance.

\begin{definition}[$\kappa$-Rate]
The $\kappa$-rate is the convergence intensity of a perpetual
financial instrument linking a price claim to an underlying
real asset. It satisfies:
\begin{enumerate}[leftmargin=*,noitemsep]
\item \textbf{Endogenous determination.} $\kappa$ is backed out
from observed market prices:
$\hat\kappa_t = (f_t/X_t)\cdot(r_b-r_a)/(1-f_t/X_t)$.
It reflects real economic convergence dynamics, not central
bank policy.
\item \textbf{Observable and transparent.} $\kappa$ is computable
from any two market prices ($f_t$ and $X_t$), publishable
on-chain, and independently verifiable.
\item \textbf{No riba content.} $\kappa$ is an event intensity,
not a reward for the time value of money. It compensates for
credit risk (sukuk), actuarial risk (takaful), or convergence
risk (perpetuals), all of which are real economic risks.
\item \textbf{Shariah-compatible.} Rate $\kappa$ can be reviewed
and validated by the Shariah board without jurisprudential
objection, as it represents neither interest on capital nor
predetermined excess.
\end{enumerate}
\end{definition}

\subsection{The $\kappa$-Yield Curve}

Just as conventional finance builds a yield curve $r(T)$ as a
function of maturity $T$, Islamic finance can build a
\textbf{$\kappa$-yield curve}: the mapping $\kappa(T)$ of
convergence intensities from perpetual sukuk of different
expected maturities.

\begin{equation}
\kappa(T) = \frac{1}{T}, \qquad T = \text{expected horizon}
\label{eq:kappa_curve}
\end{equation}
This is derived from the exponential distribution: if a sukuk is
expected to mature in $T$ years, its stopping time intensity is
$\kappa = 1/T$. The $\kappa$-yield curve is:
\begin{itemize}[leftmargin=*,noitemsep]
\item Downward sloping by construction (shorter assets have higher
$\kappa$, longer assets lower $\kappa$).
\item Grounded in real project timelines, not monetary policy
expectations.
\item A complete substitute for the conventional yield curve in
Islamic finance contexts.
\end{itemize}

\subsection{Calibration of $\kappa$}

In practice, $\kappa$ is calibrated differently for each
instrument type:
\begin{center}
\small
\begin{tabular}{@{}lll@{}}
\toprule
\textbf{Instrument} & \textbf{$\kappa$ calibration} & \textbf{Data source}\\
\midrule
Perpetual sukuk & Expected project timeline & Engineering reports\\
Takaful & Historical claim frequency & Actuarial tables\\
iCDS & Reference entity default history & Rating agencies\\
Sovereign sukuk & Project completion probability & Government data\\
\bottomrule
\end{tabular}
\end{center}

%% ============================================================
\section{Shariah Analysis}
\label{sec:shariah}
%% ============================================================

\subsection{Riba: Why $\iota = 0$ Satisfies All Schools}

The prohibition of \textit{riba} is universal across all four
Sunni schools (Hanafi, Maliki, Shafi'i, Hanbali) and the Ja'fari
(Shi'a) school \cite{Usmani2002,ElGamal2006}. The Quranic
prohibition is absolute (Al-Baqarah 2:275--279; Al-Imran 3:130).

\textit{Riba} in financial transactions takes two forms:
\begin{enumerate}[leftmargin=*,noitemsep]
\item \textit{Riba al-fadl}: excess in exchange (one party
receives more than given in a spot exchange of the same commodity).
The CEX interest floor $\iota > 0$ is precisely this: longs always
pay more than the market premium, with the excess ($\iota \cdot
f_t$) predetermined and fixed.
\item \textit{Riba al-nasi'ah}: excess for delay (predetermined
return for deferred payment). The conventional risk-free rate $r$
is this: the lender receives $r$ per unit time regardless of
economic outcome, purely for the passage of time.
\end{enumerate}

The $\kappa$-rate is neither. It is:
\begin{itemize}[leftmargin=*,noitemsep]
\item Not predetermined: $\kappa$ fluctuates with observed market
conditions and is backed out from market prices, not contractually
fixed.
\item Not a reward for time: $\kappa$ compensates for the
probability and timing of a credit/claim/convergence event
(economic risk), not for the passage of time per se.
\item Symmetric: the stopping time distribution does not
systematically benefit either party --- both carry the risk
that $\tau$ arrives sooner or later than expected.
\end{itemize}

\begin{remark}
SeekersGuidance (2025) --- the most rigorous contemporary
conservative Shariah ruling on perpetual futures --- explicitly
concedes that \textit{``removing the interest component resolves
the riba objection''}. This ruling validates our mathematical
finding: the structural design choice $\iota = 0$ (equivalently,
$r = 0$ in credit instruments) is the jurisprudentially decisive
criterion. The $\kappa$-rate produces a framework from which
the interest component has been removed by mathematical
construction, not legal re-labelling.
\end{remark}

\subsection{Gharar: Is the Stopping Time Uncertain?}

A potential objection is that the random stopping time $\tau$
introduces excessive uncertainty (\textit{gharar}) --- one party
does not know when they will receive payment.

This objection conflates two different kinds of uncertainty.
\textit{Gharar} in Islamic jurisprudence refers to uncertainty
about the \textit{existence} of the object of transaction, not
about its \textit{timing} when the timing is actuarially
determined \cite{ElGamal2006,Kamali2007}.

The stopping time $\tau \sim \mathrm{Exp}(\kappa)$ is:
\begin{enumerate}[leftmargin=*,noitemsep]
\item Precisely specified by its distribution, which is public
and on-chain verifiable.
\item Calibrated from real economic data ($\kappa$ from actuarial
tables, project timelines, or market observables).
\item Analogous to the random maturity of an ijara (lease) that
terminates when the lessee exercises an option --- a widely
accepted structure in Islamic finance \cite{AAOIFI2017}.
\end{enumerate}
The stopping time introduces actuarial uncertainty (priced by
$\kappa$), not contractual ambiguity (gharar). The two parties
know the pricing formula and the distribution of payoff times.
This is transparency, not obscurity.

\subsection{Maysir: Is This Speculation?}

\textit{Maysir} (gambling) is prohibited because it is zero-sum
and divorced from productive economic activity. The $\kappa$-rate
framework is not maysir because:
\begin{enumerate}[leftmargin=*,noitemsep]
\item Every instrument is backed by a real asset ($X_t$).
The sukuk holder's claim is on a real infrastructure project;
the takaful participant's protection covers a real agricultural
risk.
\item The premium payments $\pi_t = \kappa(f_t - X_t)\,dt$ are
transfers between economically motivated parties (capital
providers and asset operators) based on observed market
conditions, not arbitrary bets.
\item Leverage can be constrained by the $\kappa$-parameterized
framework: since $f_t = X_t$ at $\iota = 0$ (no permanent
basis), leveraged positions face immediate mark-to-market
settlement, preventing the accumulation of speculative positions
divorced from underlying values.
\end{enumerate}

\subsection{AAOIFI Compatibility}

AAOIFI Shariah Standard No.~21 (Financial Paper) and
Standard No.~17 (Investment Sukuk) govern Islamic capital market
instruments \cite{AAOIFI2017}. The $\kappa$-rate framework is
compatible because:
\begin{itemize}[leftmargin=*,noitemsep]
\item \textbf{Standard 17 (Sukuk).} The Islamic Perpetual Sukuk
satisfies the requirement that sukuk represent ownership in
an asset. The stopping time redemption mechanism is equivalent
to a put option held by the issuer --- analogous to the
``purchase undertaking'' (wa'd) structures already approved
under AAOIFI Standard 17 \cite{AdamThomas2004}.
\item \textbf{Standard 26 (Takaful).} The everlasting takaful
formula satisfies the Takaful Standard requirement for actuarial
fairness and mutual risk-sharing. The pool contributions are
actuarially determined, not based on investment returns.
\item \textbf{Standard 21 (Financial Paper).} The iCDS, as a
bilateral risk-transfer mechanism for credit risk, is analogous
to the wa'd-based credit enhancement already permitted under
Standard 21, provided the reference obligation is itself
Shariah-compliant.
\end{itemize}

%% ============================================================
\section{Discussion}
\label{sec:discussion}
%% ============================================================

\subsection{Why This Was Not Discovered Earlier}

The connection between the Ackerer perpetual contract framework
and the Duffie-Singleton credit risk model went unnoticed for
two reasons. First, the two literatures operate in different
communities (mathematical finance vs.\ fixed income) and are
rarely cited together. Second, the Ackerer framework itself is
very recent (2024), predating by only two years the present
observation. Prior to 2024, there was no mathematically rigorous
proof that $\iota = 0$ achieves no-arbitrage convergence in
perpetual contracts; the equivalence could not therefore be
established.

The deeper reason is conceptual. Conventional finance fuses two
distinct economic concepts --- the timing of a payoff event and
the time value of money --- into a single parameter $r$ (the
risk-free discount rate). This fusion appeared necessary: it
seemed that without $r$, the only way to account for the
time-dependence of a payoff was to specify it as a fixed
contractual date. The Ackerer framework dissolves this
presumption: the stopping time distribution $\mathrm{Exp}(\kappa)$
encodes timing information probabilistically, without any
predetermined interest.

\subsection{The Separation Theorem}

We can state the conceptual contribution as a general principle:

\begin{theorem}[Separation of Credit Risk from Time Preference]
\label{thm:separation}
In the Ackerer-Hugonnier-Jermann framework, the no-arbitrage
price of a financial claim $\phi(X_\tau)$ at $\iota = 0$ can
be written as:
\begin{equation}
P_t = \underbrace{\E^Q[\phi(X_\tau)]}_{\text{credit risk term}}
\quad\text{without any}\quad
\underbrace{\mathrm{e}^{-r(s-t)}}_{\text{time preference term}}
\label{eq:separation}
\end{equation}
The timing of the payoff is embedded in the distribution of
$\tau \sim \mathrm{Exp}(\kappa)$. The value of the payoff is
embedded in the expectation $\E[\phi(X_\tau)]$. Time preference
(interest) is a third, separable component that can be set to
zero without violating no-arbitrage.
\end{theorem}

This separation is the mathematically precise statement of what
Islamic finance has sought for 50 years: a proof that interest
is not necessary for no-arbitrage pricing of time-dependent
financial instruments.

\subsection{Limitations}

Several limitations qualify the current analysis.

\textbf{Constant-parameter assumption.} The closed-form results
in Sections~\ref{sec:sukuk}--\ref{sec:sovereign} assume constant
$\kappa$ and $\sigma_X$. Stochastic $\kappa$ (time-varying
hazard rates) require numerical methods; the equivalence
(Proposition~\ref{prop:equiv}) holds in the stochastic case,
but closed forms are generally unavailable.

\textbf{Recovery structure.} We assume full recovery
($\phi(X_\tau) = X_\tau$) or floor-only recovery
($\phi(X_\tau) = (K-X_\tau)^+$). Partial recovery models
(fractional loss given default) require adjusting $\phi$ accordingly.

\textbf{Liquidity.} The theory assumes liquid markets for both
$f_t$ (the perpetual instrument) and $X_t$ (the underlying asset).
For infrastructure sukuk or agricultural takaful, the underlying
may be illiquid. Liquidity adjustment to the $\kappa$-rate is
an open research question.

\textbf{Regulatory recognition.} The $\kappa$-rate framework
requires regulatory acceptance in each jurisdiction. AAOIFI
compatibility is argued above on principles grounds; formal
adoption requires engagement with the AAOIFI Standard-setting
process, which is ongoing.

%% ============================================================
\section{Conclusion}
\label{sec:conclusion}
%% ============================================================

We have established that the random time representation of
perpetual futures contracts --- a result from contemporary
mathematical finance \cite{Ackerer2024} --- is formally
isomorphic to the reduced-form credit risk pricing formula of
Duffie \& Singleton \cite{DuffieSingleton1999}. The isomorphism
maps the convergence intensity $\kappa$ to the default hazard
rate $\lambda$, and the interest parameter $\iota$ to the
risk-free discount rate $r$. Setting $\iota = r = 0$, which
Ackerer et al.\ prove achieves unique no-arbitrage pricing for
perpetual contracts, yields a complete, no-arbitrage pricing
framework for credit instruments that requires no positive
interest rate.

The implications for Islamic finance are profound. The $\kappa$-rate
provides a mathematically rigorous, Shariah-compatible pricing
foundation for: (i) Islamic Perpetual Sukuk, priced at the
expected asset value at a random redemption event; (ii) Everlasting
Takaful contracts, with premiums given by the closed-form
everlasting put formula; (iii) Islamic Credit Default Swaps,
transferring credit risk without any interest component; and (iv)
Sovereign Perpetual Sukuk for infrastructure finance, with
$\kappa$ calibrated from project timelines.

The deeper contribution is a Separation Theorem (Theorem
\ref{thm:separation}): interest is a separable, eliminable
component of no-arbitrage pricing, not an intrinsic mathematical
necessity. The stopping time distribution carries all timing
information previously attributed to the interest discount.
This is the precise mathematical statement that Islamic finance
has required but lacked: riba is not a logical prerequisite for
fair pricing of time-dependent financial instruments.

Taken together, these results suggest that the Ackerer-Hugonnier-Jermann
framework may provide the mathematical foundation for a complete
riba-free financial system --- not as a legal workaround, but as
a rigorously proved alternative architecture. Future work should
address stochastic $\kappa$ processes, liquidity premia in the
$\kappa$-yield curve, and the macroeconomic properties of a
monetary system anchored to $\kappa$ rather than $r$.

\bigskip
\noindent\textit{``And Allah has permitted trade and forbidden
\textit{riba}.''}\\
\textit{(Al-Qur'an, Al-Baqarah 2:275)}

%% ============================================================
%% REFERENCES
%% ============================================================
\bibliographystyle{plainnat}
\begin{thebibliography}{99}

\bibitem{Ackerer2024}
Ackerer, D., Hugonnier, J., \& Jermann, U.~J. (2024).
Perpetual futures pricing.
\textit{Mathematical Finance}, \textbf{34}(4), 1277--1308.
\url{https://doi.org/10.1111/mafi.12412}

\bibitem{AAOIFI2017}
AAOIFI. (2017).
\textit{Shariah Standards for Islamic Financial Institutions}
(Standard No.~17: Investment Sukuk; Standard No.~21: Financial
Paper; Standard No.~26: Islamic Insurance).
Accounting and Auditing Organization for Islamic Financial
Institutions, Manama, Bahrain.

\bibitem{AdamThomas2004}
Adam, N.~J., \& Thomas, A. (2004).
\textit{Islamic Bonds: Your Guide to Issuing, Structuring and
Investing in Sukuk}.
Euromoney Books, London.

\bibitem{Ahmed2026}
Ahmed, S., Bhuyan, R., \& Islam, R. (2026).
The interest parameter in perpetual futures: Shariah analysis
and empirical evidence from centralized and decentralized exchanges.
\textit{Working Paper}, Arcus Quant Fund.

\bibitem{BIS2024}
Bank for International Settlements. (2024).
\textit{OTC Derivatives Statistics at End-December 2023}.
BIS Quarterly Review, March 2024.
\url{https://www.bis.org/statistics/derstats.htm}

\bibitem{BlackCox1976}
Black, F., \& Cox, J.~C. (1976).
Valuing corporate securities: Some effects of bond indenture
provisions.
\textit{Journal of Finance}, \textbf{31}(2), 351--367.

\bibitem{BlackScholes1973}
Black, F., \& Scholes, M. (1973).
The pricing of options and corporate liabilities.
\textit{Journal of Political Economy}, \textbf{81}(3), 637--654.

\bibitem{CIR1985}
Cox, J.~C., Ingersoll, J.~E., \& Ross, S.~A. (1985).
A theory of the term structure of interest rates.
\textit{Econometrica}, \textbf{53}(2), 385--408.

\bibitem{DuffieLando2001}
Duffie, D., \& Lando, D. (2001).
Term structures of credit spreads with incomplete accounting
information.
\textit{Econometrica}, \textbf{69}(3), 633--664.

\bibitem{DuffieSingleton1999}
Duffie, D., \& Singleton, K.~J. (1999).
Modeling term structures of defaultable bonds.
\textit{Review of Financial Studies}, \textbf{12}(4), 687--720.

\bibitem{ElGamal2006}
El-Gamal, M.~A. (2006).
\textit{Islamic Finance: Law, Economics, and Practice}.
Cambridge University Press, Cambridge.

\bibitem{HarrisonKreps1979}
Harrison, J.~M., \& Kreps, D.~M. (1979).
Martingales and arbitrage in multiperiod securities markets.
\textit{Journal of Economic Theory}, \textbf{20}(3), 381--408.

\bibitem{Htay2012}
Htay, S.~N.~N., Arif, M., Soualhi, Y., Zaharin, H.~R., \&
Shaugee, I. (2012).
\textit{Accounting, Auditing and Governance for Takaful
Operations}.
John Wiley \& Sons, Singapore.

\bibitem{IFSB2024}
Islamic Financial Services Board. (2024).
\textit{Islamic Financial Services Industry Stability Report 2024}.
IFSB, Kuala Lumpur.

\bibitem{IIFM2023}
International Islamic Financial Market. (2023).
\textit{Sukuk Report: A Comprehensive Study of the Global Sukuk
Market} (13th ed.). IIFM, Manama.

\bibitem{IqbalMolyneux2005}
Iqbal, M., \& Molyneux, P. (2005).
\textit{Thirty Years of Islamic Banking: History, Performance,
and Prospects}.
Palgrave Macmillan, London.

\bibitem{JarrowTurnbull1995}
Jarrow, R.~A., \& Turnbull, S.~M. (1995).
Pricing derivatives on financial securities subject to credit
risk.
\textit{Journal of Finance}, \textbf{50}(1), 53--85.

\bibitem{Kamali2007}
Kamali, M.~H. (2007).
\textit{Islamic Commercial Law: An Analysis of Futures and Options}.
The Islamic Texts Society, Cambridge.

\bibitem{Khan2010}
Khan, F. (2010).
How `Islamic' is Islamic banking?
\textit{Journal of Economic Behavior \& Organization},
\textbf{76}(3), 805--820.

\bibitem{KhanAbdallah2017}
Khan, T., \& Abdallah, A. (2017).
Rethinking sukuk pricing: The ghost of LIBOR.
\textit{International Journal of Islamic and Middle Eastern
Finance and Management}, \textbf{10}(4), 492--511.

\bibitem{Lando1998}
Lando, D. (1998).
On Cox processes and credit risky securities.
\textit{Review of Derivatives Research}, \textbf{2}(2--3), 99--120.

\bibitem{Leland1994}
Leland, H.~E. (1994).
Corporate debt value, bond covenants, and optimal capital structure.
\textit{Journal of Finance}, \textbf{49}(4), 1213--1252.

\bibitem{Merton1974}
Merton, R.~C. (1974).
On the pricing of corporate debt: The risk structure of interest
rates.
\textit{Journal of Finance}, \textbf{29}(2), 449--470.

\bibitem{Usmani2002}
Usmani, M.~T. (2002).
\textit{An Introduction to Islamic Finance}.
Kluwer Law International, The Hague.

\end{thebibliography}

\end{document}
